%! AP Statistics — Unit 2, Lesson 2.3 — Slides
\documentclass[aspectratio=169, 11pt]{beamer}
\usepackage{apstats-beamer}

\begin{document}

% ── TITLE SLIDE ───────────────────────────────────────────────────────────────
{
\setbeamercolor{background canvas}{bg=navy}
\begin{frame}[plain]
  \vspace{2.8cm}\hspace{0.55cm}%
  \begin{minipage}{0.9\textwidth}
    {\color{goldacc}\bfseries\small LESSON 2.3}\\[0.5cm]
    {\color{white}\bfseries\LARGE Statistics for Two Categorical Variables}\\[1.2cm]
    {\color{skymid}\small AP Statistics~$\cdot$~Unit 2: Exploring Two-Variable Data}
  \end{minipage}
\end{frame}
}

% ── LEARNING TARGETS ─────────────────────────────────────────────────────────
\begin{frame}[t, plain]
  \navyheader{Today's Learning Targets}
  \vspace{0.3cm}
  By the end of class, I will be able to\dots
  \begin{itemize}
    \item compute the \textbf{conditional proportion} ($\hat{p}$) for each group in a two-way table.
    \item calculate and interpret the \textbf{difference in proportions} in context.
    \item calculate and interpret \textbf{relative risk} in context.
    \item compare both statistics and explain which better describes the association.
  \end{itemize}
  \vspace{0.5cm}
  \textcolor{goldacc}{\small Standards: UNC-1.Q [Skill 2.C] $\cdot$ UNC-1.R [Skill 2.D]}
\end{frame}

% ── HOOK ─────────────────────────────────────────────────────────────────────
\begin{frame}[t, plain]
  \navyheader{Hook --- Same Data, Two Stories}
  \vspace{0.3cm}
  A clinical study reports:
  \begin{itemize}
    \item Drug A group: \textbf{2\%} got the side effect.
    \item Placebo group: \textbf{1\%} got the side effect.
  \end{itemize}
  \vspace{0.5cm}
  \begin{columns}
    \column{0.48\textwidth}
      \textbf{Story 1 (Additive):}\\
      The drug group's rate is \textbf{1 percentage point higher}.\\[0.4cm]
      \textit{``Not a big deal\ldots''}
    \column{0.48\textwidth}
      \textbf{Story 2 (Multiplicative):}\\
      Drug group is \textbf{twice as likely} to get the side effect.\\[0.4cm]
      \textit{``That sounds alarming!''}
  \end{columns}
  \vspace{0.6cm}
  \textcolor{navy}{\small Both statements are mathematically correct. Today we learn the statistics behind each.}
\end{frame}

% ── DIFFERENCE IN PROPORTIONS ────────────────────────────────────────────────
\begin{frame}[t, plain]
  \navyheader{Statistic 1 --- Difference in Proportions}
  \vspace{0.3cm}
  \[
    \hat{p}_1 - \hat{p}_2 \;=\;
    \frac{\text{successes in Group 1}}{\text{Group 1 total}}
    -
    \frac{\text{successes in Group 2}}{\text{Group 2 total}}
  \]
  \vspace{0.3cm}
  \textbf{Interpretation template:}\\[0.2cm]
  \textit{``The proportion of [outcome] among [Group 1] was [value] [higher / lower] than the
  proportion among [Group 2].''}

  \vspace{0.5cm}
  \textbf{Key properties:}
  \begin{itemize}
    \item Value of \textbf{0} $\Rightarrow$ no difference, no association.
    \item Units are in proportions (or percentage points).
    \item Sign tells you direction.
  \end{itemize}
\end{frame}

% ── RELATIVE RISK ─────────────────────────────────────────────────────────────
\begin{frame}[t, plain]
  \navyheader{Statistic 2 --- Relative Risk}
  \vspace{0.3cm}
  \[
    \text{RR} \;=\; \frac{\hat{p}_1}{\hat{p}_2}
  \]
  Group 2 is the \textbf{baseline} (reference) group.

  \vspace{0.3cm}
  \textbf{Interpretation template:}\\[0.2cm]
  \textit{``[Group 1] was [RR] times as likely to [outcome] as [Group 2].''}

  \vspace{0.5cm}
  \textbf{Key properties:}
  \begin{itemize}
    \item Value of \textbf{1} $\Rightarrow$ no difference, no association.
    \item RR $> 1$: Group 1 higher. \quad RR $< 1$: Group 1 lower.
    \item \textbf{Not} a percentage-point difference --- RR $= 2$ means ``twice as likely,'' not ``2\% more.''
  \end{itemize}
\end{frame}

% ── COMPARING THE STATISTICS ─────────────────────────────────────────────────
\begin{frame}[t, plain]
  \navyheader{When Do They Tell Different Stories?}
  \vspace{0.3cm}
  \begin{columns}
    \column{0.48\textwidth}
      \textbf{Rare outcome:}\\
      Group 1: $\hat{p}_1 = 0.002$ \\
      Group 2: $\hat{p}_2 = 0.001$ \\[0.3cm]
      Difference: $+0.001$ (tiny)\\
      RR: $2.0$ (striking!)\\[0.3cm]
      \textit{RR tells a more dramatic story.}
    \column{0.48\textwidth}
      \textbf{Common outcome:}\\
      Group 1: $\hat{p}_1 = 0.60$ \\
      Group 2: $\hat{p}_2 = 0.40$ \\[0.3cm]
      Difference: $+0.20$ (meaningful!)\\
      RR: $1.5$ (moderate)\\[0.3cm]
      \textit{Absolute difference is more interpretable.}
  \end{columns}
  \vspace{0.5cm}
  \textcolor{navy}{\bfseries Takeaway:} Always report \emph{both} statistics in context.
  A small absolute difference can correspond to a large relative risk --- and vice versa.
\end{frame}

% ── CLOSING ──────────────────────────────────────────────────────────────────
{
\setbeamercolor{background canvas}{bg=navy}
\begin{frame}[plain]
  \vspace{3cm}\hspace{0.55cm}%
  \begin{minipage}{0.9\textwidth}
    {\color{goldacc}\bfseries\large Exit Ticket}\\[0.6cm]
    {\color{white} A health survey finds 45 of 150 smokers (30\%) and 18 of 180 non-smokers (10\%)
    developed a condition.}\\[0.5cm]
    {\color{skymid}\small
      (1) Compute the difference in proportions and interpret.\\
      (2) Compute the relative risk and interpret.\\
      (3) Which statistic better conveys the strength of association? Why?}
  \end{minipage}
\end{frame}
}

\end{document}
